% !TEX root = Thesis.tex

%==============================================================================
\chapter{Diffraction of a \acl{bec} with a one-dimensional Optical lattice}
\label{chap:one_dimensional_lattices}
%==============================================================================
In ultracold atoms physics, Optical lattices are defined as periodic standing wave potentials generated by interfering laser beams. Its use allows to replicate results from solid state physics, with the interaction between an ultracold atomic ensemble and an optical lattice being equivalent to the role of electrons in an atomic lattice \cite{Lewenstein2007, Bloch2008}.  However, the use of ultracold systems presents some advantages when comparing with those used in solid state physics. The main one being the capability of changing the optical lattice properties, by simply adjusting the intensity and detuning of the laser beams forming it. An additional advantage of ultracold atomic systems is the non-existence of crystal defects, which can be a big source of noise in solid states systems \cite{VanDerZiel1978}. This experiment will seek to study how an erbium \ac{bec} behaves when interacting with a one-dimensional optical lattice. In this chapter, an introduction describing the theory behind optical lattices will be shown. After this, it follows a brief description of the implemented set up to form the lattice by using the \SI{841}{\nano\meter} erbium transition (see Table \ref{tab:Transitions}).

\begin{figure}[!htbp]\centering
	\includegraphics[width=1.\columnwidth]{lattice_beams.pdf}
	\caption[Scheme of an atomic \ac{bec} interacting with a one dimensional optical lattice]{Scheme of an atomic \ac{bec} interacting with a one dimensional optical lattice. For this case, it has been decided that the optical lattice is generated by two counter-propagating laser beams with detuning $\delta$ from the erbium \SI{841}{\nano\meter} transition. The lattice potential characteristics can be tuned by simply adjusting the beams intensities and $\Delta$, representing the frequency detuning between beams.}\label{fig:lattice_beams}
\end{figure}

\section{Theoretical description of an Optical lattice}

As it can be seen in Figure \ref{fig:lattice_beams}, generating an optical lattice requires the use of two interfering laser beams, with a detuning $\delta$ from the used atomic transition (in this case the \SI{841}{\nano\meter} erbium transition). The main parameter that allows to characterize the experiment is known as the lattice depth $U_0$ and will be determined theoretically in this section by assuming some approximations in the atom behaviour. Moreover, the lattice prosperities can be changed by adjusting the optical intensity of the beams as well as the frequency detuning between beams defined here as $\Delta$. To describe theoretically the interaction of an optical lattice with an atomic ensemble, first is important to characterize accurately the interaction with just one of the optical beams forming the lattice.

\subsection{Dipole interaction of one optical beam}

Taken individually, a laser beam would interact with the atomic cloud like a typical \acf{odt}. Therefore, each beam generates a dipole potential in the atomic ensemble given by Equation \eqref{eq:interaction_potential}. However, in this case there is not a far detuning between the laser beams forming the lattice and erbium atomic transitions. Due to this, assuming the optical fields as quasi-electrostatic does not work any more. But, Equation \eqref{eq:interaction_potential} holds in any case and can be rewritten when the atoms are approximated as a two level scheme following the behaviour of a classical Lorentz oscillator \cite{Grimm2000}. For this case, it can be proved that the complex polarizability $\alpha$ has the form:

\begin{equation}\label{eq:complex_polarizability}
	\alpha = 6 \pi \epsilon_0 c^3 \frac{\Gamma/\omega_0^2}{\omega_0^2 -\omega^2 -i(\omega^3/\omega_0^2)\Gamma} 
\end{equation} 

Where $\omega$ is the optical frequency of the laser beam interacting with the atom ensemble, $\omega_0$ is the atomic transition and $\Gamma$ the on-resonance damping rate. For the classical Lorentz oscillator, this rate is equivalent to $\Gamma = \Gamma_\text{cls}$ where:

\begin{equation}\label{eq:classical_damping_rate}
	\Gamma_\text{cls} = \frac{e^2 \omega_0^2}{6\pi\epsilon_0m_ec^3}
\end{equation}

With $e$ and $m_e$ representing the electron charge and mass respectively. Even though this result is obtained when using the classical Lorentz approximation, an equivalent result can be obtained when considering the semi-classical model of the atom. In which, the atom is considered as a two-level quantum system that interacts with a classical field of light. For the semi-classical model, Equation \eqref{eq:complex_polarizability} is still valid, except for the damping rate $\Gamma$. Now, it corresponds to the spontaneous decay rate of the excited level $\Gamma_\text{smcls} \equiv \gamma$ and is given by

\begin{equation}\label{eq:semiclassical_damping_rate}
	\Gamma_\text{smcls} = \frac{\omega_0^3}{3\pi\epsilon_0 \hbar c^3} \mathopen|\bra{e}\mu\ket{g}\mathclose|^2
\end{equation}

Where $\bra{e}\mu\ket{g}$ stands for the dipole matrix element between ground and excited state. By using Equations \eqref{eq:interaction_potential} and \eqref{eq:complex_polarizability}, for the case of large enough detuning $\delta$ ( making the scattering rate much smaller than $\Gamma$) and low enough intensities $I(\vec{r})$ (avoiding saturation of the excited state), the dipole potential $U_{\text{dip}}(\vec{r})$ can be obtained as

\begin{equation}
	U_{\text{dip}}(\vec{r}) = -\frac{3\pi c^2}{2\omega_0^3} \bigg(\frac{\Gamma}{\omega_0-\omega} + \frac{\Gamma}{\omega_0+\omega}\bigg) \cdot I(\vec{r})
\end{equation} 

Being this equation valid for both the classical and semi-classical model. For most cases, the rotating wave approximation can be applied here because usually the detuning $\mathopen|\delta\mathclose| << \omega_0$. Using this approximation, the dipole potential for one laser beam interacting with an neutral atomic ensemble can be obtained as

\begin{equation}\label{eq:relation_potential_intensity}
	U_{\text{dip}}(\vec{r}) = \frac{3\pi c^2 \Gamma}{2\omega_0^3 \delta} \cdot I(\vec{r})
\end{equation} 

And by assuming the beam to be a Gaussian beam, Equations \eqref{eq:intensity_gaussian} and \eqref{eq:relation_potential_intensity} lead to the exponential behaviour already described in Equation \eqref{eq:dipole_potential}.

\subsection{One-dimensional optical lattice}

After considering only the interaction with just one beam, the optical lattice potential $U_\text{lat}(\vec{r})$, generated by two optical beams, has the same relation with the lattice intensity $I_\text{lat}(\vec{r})$ than the previous case described in Equation \eqref{eq:relation_potential_intensity}. Therefore, it can be assumed similarly that:

\begin{equation}\label{eq:relation_lattice_potential_intensity}
	U_{\text{lat}}(\vec{r}) = \frac{3\pi c^2 \Gamma}{2\omega_0^3 \delta} \cdot I_{\text{lat}}(\vec{r})
\end{equation} 

Where $I_\text{lat}(\vec{r})$ is equal to the squared absolute value of the summed electric fields $\vec{E}_1(\vec{r}, t)$ and $\vec{E}_2(\vec{r}, t)$ generated by the two beams. 

\begin{equation}\label{eq:relation_lattice_inesity}
	I_{\text{lat}}(\vec{r}) = \mathopen\big|\vec{E}_1(\vec{r}, t)+\vec{E}_2(\vec{r}, t)\mathclose\big|^2
\end{equation}

This way, assuming that both beams are linearly polarized in the arbitrary $\hat{e}_z$ direction, have the same frequency $\omega$ and are forming a perfectly parallel overlay. The electric fields have the expression:

\begin{equation}\label{eq:relation_electric_fields}
	\vec{E}_{1/2}(\vec{r}, t) = \hat{e}_z \sqrt{I_{1/2}(\vec{r})}\cdot e^{i(\pm kr -\omega t)}
\end{equation}

With $I_{1/2}(\vec{r})$ the individual intensity of the beams.


%%% Local Variables: 
%%% mode: latex
%%% TeX-master: "Thesis"
%%% End: 
